\section*{Hoofdstuk \ref{ch:g6:classification-kinship} --- Classificatie en verwantschap}

\begin{solution}{exo:g6:classification-kinship:1}
Een eigenschap die een organisme bezit, na te kijken geformuleerd: heeft
een \omterm{def:g4:classifying-animals:vertebrate}{ruggengraat}; heeft vier \omterm{def:g1:my-body:parts}{ledematen}; heeft \omterm{def:g1:animals-around-us:covering}{vacht} en zoogt zijn jongen
(of: heeft \omterm{def:g1:animals-around-us:covering}{veren}).
\end{solution}

\begin{solution}{exo:g6:classification-kinship:2}
De verzameling bezitters van elk \omterm{def:g6:classification-kinship:attribute}{kenmerk} zit volledig binnen die van een
ander: alle vachtdragers hebben vier \omterm{def:g1:my-body:parts}{ledematen}, alle viervoetige \omterm{def:g1:animals-around-us:animal}{dieren}
hebben een \omterm{def:g4:classifying-animals:vertebrate}{ruggengraat} --- nooit omgekeerd, nooit kruisend. Dozen binnen
dozen, zonder overlopers.
\end{solution}

\begin{solution}{exo:g6:classification-kinship:3}
Omdat ze die van een gemeenschappelijke voorouder geërfd hebben: een
\omterm{def:g6:classification-kinship:attribute}{kenmerk} reist een familielijn af, dus zijn de bezitters ervan de
nakomelingen van die voorouder --- een familie.
\end{solution}

\begin{solution}{exo:g6:classification-kinship:4}
De forel staat alleen voorbij de stip van de \omterm{def:g4:classifying-animals:vertebrate}{ruggengraat}. De walvis staat
voorbij \omterm{def:g4:classifying-animals:vertebrate}{ruggengraat}, vier \omterm{def:g1:my-body:parts}{ledematen} en vacht-en-melk --- alle drie.
\end{solution}

\begin{solution}{exo:g6:classification-kinship:5}
Walvis en kat: zij delen de splitsing vacht-en-melk, de jongste die
getekend is; de forel nam bij de oudste splitsing afscheid.
\end{solution}

\begin{solution}{exo:g6:classification-kinship:6}
Het is twee keer dezelfde informatie: een doos is een tak met al zijn
twijgen --- alles voorbij een splitsing erft het \omterm{def:g6:classification-kinship:attribute}{kenmerk} van die
splitsing en zit in haar doos.
\end{solution}

\begin{solution}{exo:g6:classification-kinship:7}
Vacht-en-melk wordt binnen één familielijn overgeërfd --- het markeert de
afstamming. ``Gevormd als een vis'' kan door dezelfde waterige manier van
leven op vreemden gestempeld worden, zoals de walvis laat zien: de vorm
geleend, de familie elders. De \omterm{def:g4:classifying-animals:classification}{classificatie} deelt families in, dus
vertrouwt ze het erfstuk en wantrouwt ze de geleende jas.
\end{solution}

\begin{solution}{exo:g6:classification-kinship:8}
Forel: alleen \omterm{def:g4:classifying-animals:vertebrate}{ruggengraat}. Kikker: \omterm{def:g4:classifying-animals:vertebrate}{ruggengraat}, vier \omterm{def:g1:my-body:parts}{ledematen}. Kat:
\omterm{def:g4:classifying-animals:vertebrate}{ruggengraat}, vier \omterm{def:g1:my-body:parts}{ledematen}, vacht-en-melk. Merel: \omterm{def:g4:classifying-animals:vertebrate}{ruggengraat}, vier
\omterm{def:g1:my-body:parts}{ledematen}, \omterm{def:g1:animals-around-us:covering}{veren}. Het nesten blijkt uit de vinkjes: bij elke rij zit de
\omterm{def:g4:classifying-animals:vertebrate}{ruggengraat} erbij; elke rij met een bekledingsvinkje heeft ook een vinkje
bij vier \omterm{def:g1:my-body:parts}{ledematen}.
\end{solution}

\begin{solution}{exo:g6:classification-kinship:9}
Ze passeert \omterm{def:g4:classifying-animals:vertebrate}{ruggengraat} en vier \omterm{def:g1:my-body:parts}{ledematen}, en gaat op haar eigen tak weg
met haar droge schubbige \omterm{def:g1:the-five-senses:sense}{huid} --- naast, niet voorbij, de splitsingen van
\omterm{def:g1:animals-around-us:covering}{vacht} en \omterm{def:g1:animals-around-us:covering}{veren}.
\end{solution}

\begin{solution}{exo:g6:classification-kinship:10}
Het zogen markeert de \omterm{prop:g6:classification-kinship:kinship}{verwantschap}: vacht-en-melk is het erfstuk van de
familie van de \omterm{prop:g3:sorting-living-things:groups}{zoogdieren}, dus vertakt de vleermuis samen met de kat. Het
vliegen is via de manier van leven geleend --- vleugels ontstonden apart
op de vogeltak en op de vleermuistak, zoals de stroomlijn dat deed bij
walvis en forel.
\end{solution}

\begin{solution}{exo:g6:classification-kinship:11}
Hij hoort erbij. Het vogelbekdier stamt af van de gemeenschappelijke
voorouder van de \omterm{prop:g3:sorting-living-things:groups}{zoogdieren} en draagt de erfstukken van de familie
(\omterm{def:g1:animals-around-us:covering}{vacht}, \omterm{def:g2:animals-grow-up:born}{melk}); zijn \omterm{def:g2:animals-grow-up:born}{ei} is een oud \omterm{def:g6:classification-kinship:attribute}{kenmerk} dat zijn vroege tak behield
terwijl de andere het verloren. De familie is de voorouder en
\emph{alle} nakomelingen --- zonderlingen inbegrepen; de bundel
(\cref{exo:g4:classifying-animals:11}) stemde hem er al in.
\end{solution}

\begin{solution}{exo:g6:classification-kinship:12}
Als elk paar takken bij elkaar komt, dan moeten \omterm{def:g5:biodiversity:species}{soorten} over lange tijd
veranderen --- met nieuwe \omterm{def:g6:classification-kinship:attribute}{kenmerken} die ontstaan en overgeërfd worden ---
zodat één voorouderlijke lijn zich in forel, kat en madeliefje kon
vertakken. De implicatie in één zin: alle \omterm{def:g6:exploring-our-environment:components}{levende wezens} stammen, met
verandering, van gemeenschappelijke voorouders af. Bewijs om te eisen:
sporen van de verandering --- tussenvormen, bewaarde resten in oude
gesteenten, en familiegelijkenissen precies daar waar de \omterm{def:g6:classification-kinship:tree}{stamboom} ze
voorspelt (het werk van een later jaar).
\end{solution}

\begin{solution}{pb:g6:classification-kinship:1}
\textbf{1.} Forel: \omterm{def:g4:classifying-animals:vertebrate}{ruggengraat}. Kikker: \omterm{def:g4:classifying-animals:vertebrate}{ruggengraat}, vier \omterm{def:g1:my-body:parts}{ledematen}.
Hagedis: \omterm{def:g4:classifying-animals:vertebrate}{ruggengraat}, vier \omterm{def:g1:my-body:parts}{ledematen}, droge schubbige \omterm{def:g1:the-five-senses:sense}{huid}. Duif:
\omterm{def:g4:classifying-animals:vertebrate}{ruggengraat}, vier \omterm{def:g1:my-body:parts}{ledematen}, \omterm{def:g1:animals-around-us:covering}{veren}. Muis: \omterm{def:g4:classifying-animals:vertebrate}{ruggengraat}, vier \omterm{def:g1:my-body:parts}{ledematen},
vacht-en-melk. Dolfijn: \omterm{def:g4:classifying-animals:vertebrate}{ruggengraat}, vier \omterm{def:g1:my-body:parts}{ledematen}, vacht-en-melk (de
getuigenis van het label over de \omterm{def:g2:animals-grow-up:born}{melk}; \omterm{def:g6:classification-kinship:attribute}{kenmerken} beslissen, en \omterm{def:g2:animals-grow-up:born}{melk} is
het beslissende \omterm{def:g6:classification-kinship:attribute}{kenmerk}).

\textbf{2.} De \omterm{def:g4:classifying-animals:vertebrate}{ruggengraat} --- alle zes zijn gewervelde \omterm{def:g1:animals-around-us:animal}{dieren}. De
wijdste deling markeert de oudste gemeenschappelijke voorouder: op die
diepte zijn alle zes één familie.

\textbf{3.} Kikker, hagedis, duif, muis, dolfijn. De vleugels van de duif
en de flippers van de dolfijn zijn vier \omterm{def:g1:my-body:parts}{ledematen} in werkkleding:
dezelfde geërfde \omterm{def:g1:my-body:parts}{ledematen}, gevormd door verschillende levens --- de
\omterm{def:g6:classification-kinship:tree}{stamboom} telt de erfenis, niet de kleding.

\textbf{4.} Omdat een gedeeld adres geen gedeelde familie is: ``leeft in
het water'' is een \omterm{def:g2:where-animals-live:habitat}{leefgebied}, leenbaar voor vreemden, en geen \omterm{def:g6:classification-kinship:attribute}{kenmerk}
van de afstamming van het lichaam --- de oude regel van
\cref{def:g3:sorting-living-things:criterion}, nu met haar reden erbij.

\textbf{5.} Stam $\rightarrow$ splitsing \omterm{def:g4:classifying-animals:vertebrate}{ruggengraat} (de forel gaat er
daarna uit) $\rightarrow$ splitsing vier \omterm{def:g1:my-body:parts}{ledematen} $\rightarrow$ drie
bekledingstakken: droge \omterm{def:g1:animals-around-us:covering}{schubben} (hagedis), \omterm{def:g1:animals-around-us:covering}{veren} (duif), vacht-en-melk
(muis en dolfijn samen op de jongste splitsing); de kikker gaat na de
splitsing vier \omterm{def:g1:my-body:parts}{ledematen} weg, vóór elke bekleding.

\textbf{6.} De nauwste neef van de muis in de la is de dolfijn (ze delen de
splitsing vacht-en-melk). De forel heeft ze allemaal even nauw --- ze nam
bij de eerste splitsing afscheid, dus zijn alle vijf anderen even verre
neven van haar.

\textbf{7.} Een amfibie --- viervoetig, met een kale vochtige \omterm{def:g1:the-five-senses:sense}{huid}, zonder
eigen bekledingssplitsing; zijn tak gaat weg tussen de splitsing vier
\omterm{def:g1:my-body:parts}{ledematen} en de bekledingen.

\textbf{8.} Voorbij \omterm{def:g4:classifying-animals:vertebrate}{ruggengraat} en vier \omterm{def:g1:my-body:parts}{ledematen}, de tak met de droge
\omterm{def:g1:animals-around-us:covering}{schubben} op: een reptiel --- ongezien ingedeeld, want \omterm{def:g6:classification-kinship:attribute}{kenmerken} beslissen,
geen labels.

\textbf{9.} Dolfijn en muis zitten dichter bij elkaar: vacht-en-melk is
een geërfde splitsing. De stroomlijn is het uniform van het water, dat
niet-verwante zwemmers gelijkelijk uitgereikt wordt
(\cref{exo:g4:adapted-to-their-world:11}); meegeteld zou ze geleende
kleding als familie archiveren en het nesten kapotmaken.

\textbf{10.} Eén gemeenschappelijke voorouderlijke \omterm{def:g5:biodiversity:species}{soort} --- een vroeg
viervoetig \omterm{def:g4:classifying-animals:vertebrate}{gewerveld dier} --- waarvan kikker, hagedis, duif, muis en
dolfijn allemaal afstammen, elk met zijn vier \omterm{def:g1:my-body:parts}{ledematen} in andere
kleding.

\textbf{11.} Dat alle zes --- en de klas die ze tekent --- tot één
vertakkende familie behoren, met verandering afstammend van
gemeenschappelijke voorouders, met de mensen tussen de rest op de tak van
vacht-en-melk.

\textbf{12.} Het nesten is een feit over de levende wereld --- \omterm{def:g6:classification-kinship:attribute}{kenmerken}
vallen werkelijk in dozen-binnen-dozen, en geen enkele archiefbeambte
heeft dat verordonneerd. De \omterm{def:g6:classification-kinship:tree}{stamboom} verklaart dat feit: overerving van
gemeenschappelijke voorouders \emph{moet} geneste \omterm{def:g6:classification-kinship:attribute}{kenmerken} opleveren,
terwijl louter archiveren de gewilligheid van de wereld een wonder zou
laten.
\end{solution}
